\chapter{Design: SOLID, OOP and patterns in this code}

Sir may ask "where did you use SOLID?" or "which design pattern is this?". This chapter answers with the exact code.
Each name in \cc{this font} is a function or class in this book; the purple \textbf{Design} boxes under the code repeat the point where it happens.

\begin{keyidea}[The shape of the program]
Three layers, each talking only to the one below it:
\textbf{window} (\cc{ui/}) \textrightarrow{} \textbf{controller} (\cc{project.py}) \textrightarrow{} \textbf{modules} (\cc{ingest}, \cc{renderer}, \cc{checks}, \cc{fit}\dots) \textrightarrow{} \textbf{outside programs} (pandoc, Tectonic, Ghostscript, Tesseract, through \cc{process.py}).
The window never reads a paper or a PDF itself, and the modules never touch a widget. This is \emph{separation of concerns}, and most of SOLID follows from it.
\end{keyidea}

\section{SOLID, one principle at a time}

\subsection*{S: Single Responsibility (one reason to change)}
A module or class should do one job, so one kind of change touches one place.
\begin{itemize}
\item \hyperref[u:pandoc:_run]{\cc{pandoc.\_run}} is the only code that starts pandoc; \hyperref[u:compile:compile_pdf]{\cc{compile\_pdf}} the only code that runs Tectonic; \hyperref[u:fix:embed_fonts]{\cc{embed\_fonts}} the only code that runs Ghostscript; \hyperref[u:process:run_tool]{\cc{process.run\_tool}} the only code that calls \cc{subprocess}.
\item \hyperref[u:checks:PdfChecks]{\cc{PdfChecks}} only inspects a PDF. \cc{tests/test\_layering.py} fails if \cc{checks.py} ever imports \cc{subprocess}, \cc{tkinter} or network modules: the principle is enforced by a test.
\item \hyperref[m:analyze]{\cc{analyze.py}} only judges the manuscript; \hyperref[m:renderer]{\cc{renderer.py}} only writes LaTeX; \hyperref[m:ui/style]{\cc{ui/style.py}} only holds colours.
\item Each editor class edits one kind of block (\hyperref[m:ui/editors]{\cc{ui/editors.py}}).
\end{itemize}

\subsection*{O: Open/Closed (add without editing what works)}
Code should be open to extension but closed to modification.
\begin{itemize}
\item \textbf{New block kind:} write a new \cc{Editor} subclass and add one line to \hyperref[u:ui/editors:editor_for]{\cc{editor\_for}}. No existing editor changes.
\item \textbf{New conference limits:} a new YAML file read by \hyperref[u:models:Rules.load]{\cc{Rules.load}}; no code changes.
\item \textbf{New paper template:} a new \cc{.tex.j2} file; \hyperref[u:renderer:render]{\cc{render}} takes the template name as a parameter.
\item \textbf{New AI evidence layer:} one new function and one line in \hyperref[u:images:inspect_file]{\cc{inspect\_file}}.
\item \textbf{New PDF check:} one new method in \cc{PdfChecks} and its name in the tuple in \hyperref[u:checks:PdfChecks.run]{\cc{PdfChecks.run}} (that tuple is the one line that does change).
\end{itemize}

\subsection*{L: Liskov Substitution (a subclass can stand in for its parent)}
\begin{itemize}
\item The window stores \emph{an} editor in \cc{self.editor} and only ever calls \cc{save()} and \cc{destroy()} on it (\hyperref[u:ui/window:MainWindow.save_editor]{\cc{save\_editor}}, \hyperref[u:ui/window:MainWindow.make_editor]{\cc{make\_editor}}). Every subclass of \hyperref[u:ui/editors:Editor]{\cc{Editor}} keeps the same promise: \cc{save()} writes through the Project and returns a message string. So any editor can stand in for any other.
\item Our exceptions (\cc{PandocNotFound}, \cc{TectonicNotFound}, \cc{FitFailed}\dots) subclass \cc{Exception}, so code that catches \cc{Exception} handles them too (\hyperref[u:ui/window:MainWindow.run]{\cc{MainWindow.run}} does).
\end{itemize}

\subsection*{I: Interface Segregation (do not depend on what you do not use)}
Python has no \cc{interface} keyword; the idea is to give each user only the small piece it needs.
\begin{itemize}
\item \hyperref[u:ui/outline:Outline.__init__]{\cc{Outline}} is given one function, \cc{on\_select}; \hyperref[u:ui/panels:Panels.__init__]{\cc{Panels}} one function, \cc{on\_jump}. They do not receive the whole window.
\item \hyperref[u:checks:check_fonts]{\cc{check\_fonts}}, \cc{check\_page\_count} and the other one-check functions let a caller use one check without building a \cc{PdfChecks} itself.
\item The window uses a handful of Project methods; it never sees ingest's or the fitter's internals.
\end{itemize}

\subsection*{D: Dependency Inversion (depend on a small promise, not on details)}
\begin{itemize}
\item The outline and the tabs depend on a \emph{callback} (any function taking a block id), not on \cc{MainWindow}. That is why they can be tested and reused alone.
\item \cc{MainWindow} depends on \hyperref[u:project:Project]{\cc{Project}}, not on the modules behind it.
\item Every module depends on \hyperref[u:process:find_tool]{\cc{find\_tool}} and \cc{run\_tool}, not on how each operating system finds and runs programs.
\item The rules are passed in (\cc{build(rules)}, \cc{fit(rules)}), not hard-coded; \cc{choose\_fit} is stored in a variable so tests can inject a choice (dependency injection, \hyperref[u:ui/window:MainWindow.__init__]{\cc{MainWindow.\_\_init\_\_}}).
\end{itemize}

\begin{trap}[Be honest about where SOLID is not followed]
\textbf{\cc{MainWindow}} lays out the window \emph{and} handles every action: more than one responsibility. It stays readable because each action is a short method that hands the work to \cc{Project}.\par
\textbf{\cc{ingest.py}} finds the title block and tags the body in one file; both walk the same pandoc blocks, but it could be split.\par
\textbf{Editors receive the whole window} (to call \cc{save\_editor}, \cc{move}\dots), so they depend on more than they strictly need.\par
\textbf{\cc{Project} creates concrete modules directly} instead of receiving them; for a project this size that is simpler, and still testable.\par
Saying this calmly, with the improvement you would make, is a strong viva answer.
\end{trap}

\section{The four ideas of object-oriented programming}

\begin{tabularx}{\linewidth}{>{\bfseries}l X}
\toprule
Encapsulation & Data and the code that changes it live together. Only Project's methods change a manuscript, and each saves it (\hyperref[u:project:Project.set_caption]{\cc{set\_caption}}). Helpers starting with \cc{\_} (\cc{Project.\_block}, \cc{pandoc.\_run}) are private by convention.\\
Abstraction & A simple call hides many steps: \hyperref[u:project:Project.build]{\cc{Project.build()}} renders, instruments, compiles, measures, strips markers and checks.\\
Inheritance & \cc{FrontEditor}, \cc{TextEditor}\dots{} inherit from \cc{Editor}, which inherits from \cc{ttk.Frame}; our exceptions inherit from \cc{Exception}.\\
Polymorphism & Each editor's \cc{save()} does something different under one name; \hyperref[u:models:Author.__eq__]{\cc{Author.\_\_eq\_\_}} overloads \cc{==}; \cc{@property} makes \cc{Measured.width} read like a field.\\
\bottomrule
\end{tabularx}

\section{Design patterns you can point at}

\begin{longtable}{>{\bfseries\raggedright\arraybackslash}p{0.2\linewidth} >{\raggedright\arraybackslash}p{0.74\linewidth}}
\toprule
Pattern & Where, and why\\
\midrule\endhead
Facade & \hyperref[u:project:Project]{\cc{Project}}: one simple front for many modules, so the window has a small set of calls to make.\\
Factory & \hyperref[u:ui/editors:editor_for]{\cc{editor\_for}} picks and creates the right editor class. \cc{Manuscript.load}, \cc{Rules.load} and \cc{Project.open} are factory methods (classmethods that build an object).\\
Strategy & The fit \cc{policy} ("readable" or "limit") in \hyperref[u:fit:plan]{\cc{plan}}; the replaceable \cc{choose\_fit} function in the window.\\
Observer / callback & Tk calls our functions on events (\cc{command=}, \cc{bind}); \cc{Outline} calls \cc{on\_select}, \cc{Panels} calls \cc{on\_jump}.\\
Data class / DTO & \hyperref[m:models]{\cc{models.py}}, \cc{Block}, \cc{Finding}: plain carriers of data between layers.\\
Future / producer-consumer & \hyperref[u:ui/window:MainWindow.run]{\cc{MainWindow.run}}: the worker thread produces a result, the main thread consumes it.\\
Debounce & \hyperref[u:ui/window:MainWindow.schedule_build]{\cc{schedule\_build}}: many edits, one build.\\
Cache / lazy initialisation & \hyperref[u:measure:ocr_text_per_width]{\cc{ocr\_text\_per\_width}} caches OCR by file and time; \hyperref[u:images:ai_score]{\cc{ai\_score}} loads the model once.\\
Batching & \hyperref[u:pandoc:blocks_to]{\cc{blocks\_to}}: many conversions in one pandoc run.\\
Template method (light) & \cc{Editor.\_\_init\_\_} draws the shared heading and problems; each subclass fills in its own fields and \cc{save()}.\\
\bottomrule
\end{longtable}

\section{Algorithms and data structures in the code}

\begin{tabularx}{\linewidth}{>{\bfseries}l X}
\toprule
Stack & Matching \cc{\textbackslash begin\{figure\}} with \cc{\textbackslash end\{figure\}}: \hyperref[u:figures:_figure_envs]{\cc{\_figure\_envs}}.\\
Counter for nesting & Matching braces: \hyperref[u:texsource:_group]{\cc{texsource.\_group}}, \cc{figures.\_split\_options}.\\
Recursion & Walking the pandoc tree (\cc{stringify}, \cc{\_images}), unwrapping Divs (\cc{\_flatten}), following \cc{\textbackslash input} (\cc{full\_source}).\\
Binary search & The page-limit search in \hyperref[u:fit:fit_paper]{\cc{fit\_paper}}: 5 compiles instead of dozens.\\
Hashing & dHash and Hamming distance for duplicate images: \hyperref[u:images:dhash]{\cc{dhash}}, \cc{same\_picture}.\\
Sets & Cited keys versus reference keys (\cc{analyze}); set difference finds the missing ones.\\
Median & Robust text size, ignoring odd values: \hyperref[u:measure:small_text]{\cc{small\_text}}.\\
\bottomrule
\end{tabularx}

\section{Other principles sir may name}

\begin{itemize}
\item \textbf{DRY (do not repeat yourself):} one \cc{Finding} class for every manuscript problem; one \cc{ProblemList} for three tabs; colours only in \cc{style.py}; \cc{Project.export} reuses \cc{build}.
\item \textbf{KISS (keep it simple):} plain classes in \cc{models.py}; Tkinter instead of a large framework; one worker thread instead of many.
\item \textbf{YAGNI (you are not going to need it):} only the IEEE template is built; no plugin system for templates until a second one is needed.
\item \textbf{Fail fast:} \cc{Project.\_block} raises \cc{KeyError}; \cc{Rules.load} rejects unknown keys; \cc{compile\_pdf} deletes the old PDF so failure is visible.
\item \textbf{Defensive programming:} zip-slip check in \cc{\_unzip}; \cc{yaml.safe\_load}; \cc{-dSAFER} for Ghostscript; a list of arguments instead of \cc{shell=True}.
\item \textbf{Testability:} the window logic lives in \cc{Project}, so most tests run without a window; the window's dialogs are replaceable (\cc{choose\_fit}).
\end{itemize}

\begin{cheat}{Design questions, with answers}
\textbf{Where is Single Responsibility?} Each outside program has exactly one function that runs it; \cc{checks.py} only reads PDFs, and a test enforces it.\par
\textbf{Show me Open/Closed.} \cc{editor\_for}: a new block kind is a new subclass plus one line. Rules in YAML: new limits without new code.\par
\textbf{Show me Liskov.} The window calls \cc{save()} on whatever editor it holds; every subclass returns a message string, so any one can replace another.\par
\textbf{Show me Dependency Inversion.} \cc{Outline(parent, on\_select)}: it depends on a function, not on the window. The window depends on \cc{Project}, not on the modules.\par
\textbf{Which pattern is Project?} A Facade: one simple interface over many modules.\par
\textbf{Why not more classes and interfaces?} KISS and YAGNI: for one template and one window, extra layers would make the code harder to explain without adding anything.
\end{cheat}
